\subsection{Area of a Triangle}



\begin{definition}
	The area of the rectangle $ACBD$ shown in Fig. \ref{ch2_sq_ar} is defined as $ab$. Note that all the angles in the rectangles are $90^{\degree}$
	\label{ch2_sq_ar}
\end{definition}
\begin{figure}[!ht]
	\begin{center}
		
		%\includegraphics[width=\columnwidth]{./figs/ch2_sq_ar}
		%\vspace*{-10cm}
		\resizebox{\columnwidth}{!}{\begin{tikzpicture}
[scale=2,>=stealth,point/.style={draw,circle,fill = black,inner sep=0.5pt},]

\node (C) at (0, 0)[point,label=below left:$C$] {};
\node (A) at (0, 3)[point,label=above left:$A$]{};
\node (D) at (3, 3)[point,label=above right:$D$]{};
\node (B) at (3, 0)[point,label=below right:$B$]{};

\draw (C)--(A);
\draw (A)--(D);
\draw (D)--(B);
\draw (B)--(C);
\draw (B)--(A);
\tkzMarkRightAngle[size=.2](B,C,A)
\tkzMarkRightAngle[size=.2](A,D,B)

\node [below] at (1.5,0) {$a$};
\node [above] at (1.5,3) {$a$};
\node [left] at (0,1.5) {$b$};
\node [right] at (3,1.5) {$b$};

\end{tikzpicture}}
	\end{center}
	\caption{Area of a Right Triangle}
	\label{ch2_sq_ar}	
\end{figure}

\begin{definition}
	The area of the two triangles constituting the rectangle is the same.
	\label{ch2_triang_eq}
\end{definition}
\begin{definition}
	The area of the rectangle is the sum of the areas of the two triangles inside.
	\label{ch2_triang_sum}
\end{definition}

\begin{problem}
	Show that the area of $\Delta ABC$ is $\frac{ab}{2}$
\end{problem}
\begin{figure}[!ht]
	\begin{center}
		
		%\includegraphics[width=\columnwidth]{./figs/ch2_triang_ar}
		%\vspace*{-10cm}
		\resizebox{\columnwidth}{!}{\begin{tikzpicture}
[scale=2,>=stealth,point/.style={draw,circle,fill = black,inner sep=0.5pt},]

\node (D) at (0, 0)[point,label=below :$D$] {};
\node (A) at (0, 3)[point,label=above :$A$]{};
\node (B) at (-3, 0)[point,label=below left:$B$]{};
\node (C) at (3, 0)[point,label=below right:$C$]{};

\draw (D)--(B);
\draw (B)--(A);
\draw (A)--(C);
\draw (C)--(D);
\draw (D)--(A);

\tkzMarkRightAngle[size=.2](A,D,C)

\node [below] at (0,-0.3) {$a$};
\node [below] at (-1.5,0) {$x$};
\node [below] at (1.5, 0) {$y$};
\node [below] at (0.1,1.5) {$h$};
\node [above] at (-1.5,1.5){$c$};
\node [above] at (1.5,1.5){$b$};

\end{tikzpicture}}
	\end{center}
	\caption{Area of a Triangle}
	\label{ch2_triang_ar}	
\end{figure}

\proof From\eqref{ch2_triang_sum},
\begin{equation}
\label{ch2_triang_ar_1}
ar\brak{ABCD} = ar\brak{ACB} + ar\brak{ADB}
\end{equation}
Also from \eqref{ch2_triang_eq},
\begin{equation}
\label{ch2_triang_ar_2}
ar\brak{ACB} = ar\brak{ADB}
\end{equation}
From \eqref{ch2_triang_ar_1} and \eqref{ch2_triang_ar_2},
\begin{align}
2ar\brak{ACB} &= ar\brak{ABCD} = ab \brak{\text{from} \quad \eqref{ch2_sq_ar}}
\\
\Rightarrow ar\brak{ACB} &= \frac{ab}{2}
\end{align}

\begin{problem}
	Show that the area of $\Delta ABC$ in Fig. 	\ref{ch2_triang_ar}	is $\frac{1}{2}ah$.
\end{problem}

\proof In Fig. \ref{ch2_triang_ar},
\begin{align}
ar\brak{\Delta ADC} &= \frac{1}{2}hy \\
ar\brak{\Delta ADB} &= \frac{1}{2}hx 
\end{align}
Thus,
\begin{align}
ar\brak{\Delta ABC} &= ar\brak{\Delta ADC} + ar\brak{\Delta ADB} \\
&= \frac{1}{2}hy + \frac{1}{2}hx = \frac{1}{2}h\brak{x+y} \\
&= \frac{1}{2}ah
\end{align}
\begin{problem}
	Show that the area of $\Delta ABC$ in Fig. 	\ref{ch2_triang_ar}	is $\frac{1}{2}ab \sin C$.
\end{problem}
\proof We have
%
\begin{equation}
ar\brak{\Delta ABC} = \frac{1}{2}ah = \frac{1}{2}ab\sin C \quad \brak{\because \quad h = b \sin C}.
\end{equation}
%
\begin{problem}
	Show that 
	\begin{equation}
	\frac{\sin A}{a} = \frac{\sin B}{b} = \frac{\sin C}{c}
	\end{equation}
\end{problem}
\proof Fig. \ref{ch2_triang_ar} can be suitably modified to obtain 
\begin{equation}
ar\brak{\Delta ABC} = \frac{1}{2}ab\sin C = \frac{1}{s}bc\sin A = \frac{1}{2}ca\sin B
\end{equation}
Dividing the above by $abc$, we obtain
	\begin{equation}
	\frac{\sin A}{a} = \frac{\sin B}{b} = \frac{\sin C}{c}
	\end{equation}
This is known as the sine formula.	
%
\subsection{Median}
\begin{definition}
	The line AD in Fig. \ref{ch2_median_def} that divides the side $a$ in two equal halfs is known as the median.
\end{definition}
\begin{figure}[!ht]
	\begin{center}
		
		%\includegraphics[width=\columnwidth]{./figs/ch2_median_def}
		%\vspace*{-10cm}
		\resizebox{\columnwidth}{!}{\begin{tikzpicture}
[scale=2,>=stealth,point/.style={draw,circle,fill = black,inner sep=0.5pt},]

\node (D) at (0, 0)[point,label=below :$D$] {};
\node (A) at (0, 3)[point,label=above :$A$]{};
\node (B) at (-3, 0)[point,label=below left:$B$]{};
\node (C) at (3, 0)[point,label=below right:$C$]{};

\draw (D)--(B);
\draw (B)--(A);
\draw (A)--(C);
\draw (C)--(D);
\draw (D)--(A);

\node [below] at (0,-0.3) {$a$};
\node [below] at (-1.5,0) {$a/2$};
\node [below] at (1.5, 0) {$a/2$};
\node [above] at (-1.5,1.5){$c$};
\node [above] at (1.5,1.5){$b$};

\end{tikzpicture}}
	\end{center}
	\caption{Median of a Triangle}
	\label{ch2_median_def}	
\end{figure}

\begin{problem}
	Show that the median $AD$ in Fig. \ref{ch2_median_def} divides  $\Delta ABC$ into triangles $ADB$ and $ADC$ that have equal area.
\end{problem}
\proof We have
%
\begin{align}
ar \brak{\Delta ADB} &= \frac{1}{2}\frac{a}{2}c \sin B =  \frac{1}{4}ac \sin B \\
ar \brak{\Delta ADC} &= \frac{1}{2}\frac{a}{2}b \sin C =  \frac{1}{4}ab \sin C 
\end{align}
%
Using the sine formula, $b \sin C = c \sin B$,
\begin{equation}
ar \brak{\Delta ADB} = ar \brak{\Delta ADC}
\end{equation}
\begin{problem}
	$BE$ and $CF$ are the medians in Fig. \ref{ch2_median_ratio}.  Show that
	\begin{equation}
	ar \brak{\Delta BFC} = ar \brak{\Delta BEC}
	\end{equation} 
	\label{ch2_median_eq_tri}
\end{problem}
\proof Since $BE$ and $CF$ are the medians, 

\begin{align}
ar \brak{\Delta BFC} &= \frac{1}{2} ar \brak{\Delta ABC} \\
ar \brak{\Delta BEC} &= \frac{1}{2} ar \brak{\Delta ABC} 
\end{align}
From the above, we infer that
%
\begin{equation}
ar \brak{\Delta BFC} = ar \brak{\Delta BEC}
\end{equation}


\begin{figure}[!ht]
	\begin{center}
		
		%\includegraphics[width=\columnwidth]{./figs/ch2_median_ratio}
		%\vspace*{-10cm}
		\resizebox{\columnwidth}{!}{\begin{tikzpicture}
[scale=2,>=stealth,point/.style={draw,circle,fill = black,inner sep=0.5pt},]

\node (E) at (1.5, 1.5)[point,label=above :$E$] {};
\node (F) at (-1.5, 1.5)[point,label=above :$F$] {};
\node (A) at (0, 3)[point,label=above :$A$]{};
\node (B) at (-3, 0)[point,label=below left:$B$]{};
\node (C) at (3, 0)[point,label=below right:$C$]{};
\node (a) at (0,0)[point,label=below :$a$] {};
\node (O) at (0,1)[point,label=above :$O$] {};

\draw (a)--(B);
\draw (B)--(A);
\draw (A)--(C);
\draw (C)--(a);
\draw (B)--(E);
\draw (C)--(F);
\node [above] at (-1.7,1.7) {$c$};
\node [above] at (1.7,1.7) {$b$};
\node [above] at (-2.5,.75) {$c/2$};
\node [above] at (-1,2.1) {$c/2$};
\node [above] at (2.5,.75) {$b/2$};
\node [above] at (1,2.1) {$b/2$};
%\node [above] at (1,1.3) {$p$};
%\node [above] at (-1,1.3) {$q$};

\tkzMarkAngle[size=.3](F,O,B);
\tkzMarkAngle[size=.3](C,O,E);
\draw (-0.5,1) node{$\theta$};
\draw (0.5,1) node{$\theta$};

\end{tikzpicture}}
	\end{center}
	\caption{$O$ is the Intersection of Two Medians}
	\label{ch2_median_ratio}	
\end{figure}
%
\begin{problem}
	The medians $BE$ and $CF$ in Fig. \ref{ch2_median_ratio} meet at point $O$.  Show that
	\begin{equation}
	\frac{OB}{OE} = \frac{OC}{OF} 
	\end{equation} 
\end{problem}
\proof From Problem \ref{ch2_median_eq_tri},

\begin{equation}
ar \brak{\Delta BFC} = ar \brak{\Delta BEC}
\end{equation}
%
\begin{multline}
\Rightarrow ar \brak{\Delta BOF} + ar \brak{\Delta BOC} \\
= ar \brak{\Delta BOC} + ar \brak{\Delta COE} 
\end{multline}
resulting in
%
\begin{equation}
ar \brak{\Delta BOF} 
=  ar \brak{\Delta COE} 
\end{equation}
%
Using the sine formula for area of a triangle, the above equation can be expressed as
%
\begin{align}
\frac{1}{2}OB\,OF \sin \theta &= \frac{1}{2}OC\,OE \sin \theta \\
\Rightarrow 	\frac{OB}{OE} = \frac{OC}{OF} 
\end{align}

\begin{definition}
	We know that the median of a triangle  divides it into two triangles with equal area. Using this result along with the sine formula for the area of a triangle in Fig. \ref{ch2_supp_sin},
\begin{figure}[!ht]
	\begin{center}
		
		%\includegraphics[width=\columnwidth]{./figs/ch2_supp_sin}
		%\vspace*{-10cm}
		\resizebox{\columnwidth}{!}{\begin{tikzpicture}
[scale=2,>=stealth,point/.style={draw,circle,fill = black,inner sep=0.5pt},]



\node (A) at (0, 3)[point,label=above :$A$]{};
\node (B) at (-3, 0)[point,label=below left:$B$]{};
\node (C) at (3, 0)[point,label=below right:$C$]{};
\node (D) at (0,0)[point,label=below :$D$] {};
%\node (a) at (0,-0.2) [point,label=below :$a$] {};


\draw (B)--(A);
\draw (A)--(C);
\draw (C)--(B);
\draw (A)--(D);

\node [above] at (-1.5,-.7) {$a/2$};
\node [above] at (1.5,-.7) {$a/2$};
\node [above] at (-1.7,1.7) {$c$};
\node [above] at (1.7,1.7) {$b$}; 
\node [below] at (0,-0.3) {$a$}; 

\tkzMarkAngle[size=.3](A,D,B);
\tkzMarkAngle[size=.4](C,D,A);
\draw (-.5,0.1) node{$\theta$};
\draw (.7,0.1) node{180 - $\theta$};

\end{tikzpicture}}
	\end{center}
	\caption{$\sin \theta = \sin \brak{180^{\degree}-\theta}$}
	\label{ch2_supp_sin}	
\end{figure}

\begin{align}
\frac{1}{2}\frac{a}{2}AD \sin \theta &= \frac{1}{2}\frac{a}{2}AD \sin \brak{180^{\degree}-\theta} \\
\Rightarrow \sin \theta &= \sin \brak{180^{\degree} - \theta}.
\label{sin_supp}
\end{align}
Note that our geometric definition of $\sin \theta$ holds only for $\theta < 90^{\degree}$.  \eqref{sin_supp} allows us to extend this definition for $\angle ADC > 90^{\degree}$.
\end{definition}

\subsection{Similar Triangles}
\begin{problem}
In Fig. \ref{ch2_cosine_formula}, show that
%
\begin{equation}
\cos A = \frac{b^2+c^2-a^2}{2bc}
\end{equation}
%
\
\end{problem}
\begin{figure}[!ht]
	\begin{center}
		
		%\includegraphics[width=\columnwidth]{./figs/ch2_triang_ar}
		%\vspace*{-10cm}
		\resizebox{\columnwidth}{!}{\begin{tikzpicture}
[scale=2,>=stealth,point/.style={draw,circle,fill = black,inner sep=0.5pt},]

\node (D) at (0, 0)[point,label=below :$D$] {};
\node (A) at (0, 3)[point,label=above :$A$]{};
\node (B) at (-3, 0)[point,label=below left:$B$]{};
\node (C) at (3, 0)[point,label=below right:$C$]{};

\draw (D)--(B);
\draw (B)--(A);
\draw (A)--(C);
\draw (C)--(D);
\draw (D)--(A);

\tkzMarkRightAngle[size=.2](A,D,C)

\node [below] at (0,-0.3) {$a$};
\node [below] at (-1.5,0) {$x$};
\node [below] at (1.5, 0) {$y$};
\node [below] at (0.1,1.5) {$h$};
\node [above] at (-1.5,1.5){$c$};
\node [above] at (1.5,1.5){$b$};

\end{tikzpicture}}
	\end{center}
	\caption{The cosine formula}
	\label{ch2_cosine_formula}	
\end{figure}

\proof From the figure, the first of the following equations
%
\begin{align}
a &= b \cos C + c \cos B \\
b &= c \cos A + a \cos C \\
c &= b \cos A + a \cos B
\end{align}
%
is obvious and the other two can be similarly obtained.  The above equations can be expressed in matrix form as
%
\begin{equation}
\begin{pmatrix}
0 & c & b \\
c & 0 & a \\
b & a & 0
\end{pmatrix}
\begin{pmatrix}
\cos A \\
\cos B \\
\cos C
\end{pmatrix}
= 
\begin{pmatrix}
a\\
b\\
c
\end{pmatrix}
\end{equation}
%
Using the properties of determinants,
%
\begin{align}
\cos A &= \frac{
\begin{vmatrix}
a & c & b \\
b & 0 & a \\
c & a & 0
\end{vmatrix}
	}
	{
\begin{vmatrix}
0 & c & b \\
c & 0 & a \\
b & a & 0
\end{vmatrix}
	}
	&=\frac{ab^2 + ac^2 - a^3}{abc + abc} = \frac{b^2 + c^2 - a^2}{2abc}
\end{align}
%
%
\begin{problem}
	In Fig. \ref{ch2_sim_triang}, show that $EF = \frac{a}{2}$.  
\end{problem}
%
\begin{figure}[!ht]
	\begin{center}
		
		%\includegraphics[width=\columnwidth]{./figs/ch2_median_ratio_val}
		%\vspace*{-10cm}
		\resizebox{\columnwidth}{!}{\begin{tikzpicture}
[scale=2,>=stealth,point/.style={draw,circle,fill = black,inner sep=0.5pt},]

\node (E) at (1.5, 1.5)[point,label=above :$E$] {};
\node (F) at (-1.5, 1.5)[point,label=above :$F$] {};
\node (A) at (0, 3)[point,label=above :$A$]{};
\node (B) at (-3, 0)[point,label=below left:$B$]{};
\node (C) at (3, 0)[point,label=below right:$C$]{};
\node (a) at (0,0)[point,label=below :$a$] {};
\node (O) at (0,1)[point,label=below :$O$] {};

\draw (a)--(B);
\draw (B)--(A);
\draw (A)--(C);
\draw (C)--(a);
\draw (B)--(E);
\draw (C)--(F);
\draw (E)--(F);
\node [above] at (-1.7,1.7) {$c$};
\node [above] at (1.7,1.7) {$b$};
\node [above] at (-2.5,.75) {$c/2$};
\node [above] at (-1,2.1) {$c/2$};
\node [above] at (2.5,.75) {$b/2$};
\node [above] at (1,2.1) {$b/2$};
\node [below] at (1,1.3) {$p$};
\node [below] at (-1,1.3) {$q$};
\node [above] at (-.7,.5) {$k_{p}$};
\node [above] at (.7,.5) {$k_{q}$};

\tkzMarkAngle[size=.3](F,O,B);
\tkzMarkAngle[size=.3](C,O,E);
\draw (-0.5,1) node{$\theta$};
\draw (0.5,1) node{$\theta$};

\end{tikzpicture}}
	\end{center}
	\caption{Similar Triangles}
	\label{ch2_sim_triang}	
\end{figure}

\proof Using the cosine formula for $\Delta AEF$,
%
\begin{align}
EF^2 &= \brak{\frac{b}{2}}^2 + \brak{\frac{c}{2}}^2 - 2 \brak{\frac{b}{2}}\brak{\frac{c}{2}} \cos A \\
&= \frac{b^2 + c^2 - 2bc \cos A}{4} \\
&= \frac{a^2}{4} \\
\Rightarrow EF &= \frac{a}{2}
\end{align}
%

\begin{definition}
	The ratio of sides of triangles AEF and ABC is the same.  Such triangles are known as similar triangles.
\end{definition}
\begin{problem}
	Show that similar triangles have the same angles.
\end{problem}
\proof Use cosine formula and the proof is trivial.
\begin{problem}
	Show that in Fig. \ref{ch2_sim_triang}, $EF || BC$.
\end{problem}
\proof Since $\Delta AEF \sim \Delta ABC$, $\angle AEF = \angle ACB$.  Hence the line $EF||BC$
%
\begin{problem}
	Show that $\triangle OEF \sim \triangle OEC$.
\end{problem}
\begin{problem}
%	Using Fig. \ref{ch2_median_ratio_val}, 
Show that
	\begin{align}
	\frac{OB}{OE} = \frac{OC}{OF} = 2
	\end{align}
\end{problem}
%\begin{problem}
%	In Fig. \ref{ch2_median_ratio_val}, show that
%	\begin{equation}
%	ar\brak{\Delta AFE} = 	\frac{1}{4} ar\brak{\Delta ABC}
%	\end{equation}
%	\label{ch2_small_triang}
%\end{problem}
%\proof We have
%%
%\begin{align}
%ar\brak{\Delta AFE} &= \frac{1}{2} \frac{b}{2}\frac{c}{2}\sin A = \frac{1}{4}.\frac{1}{2}bc \sin A \\
%&=\frac{1}{4}ar\brak{\Delta ABC}
%\end{align}
%
		%\includegraphics[width=\columnwidth]{./figs/ch2_median_ratio_val}
		%\vspace*{-10cm}

%
%\begin{figure}[!ht]
%	\begin{center}
%		
%		\resizebox{\columnwidth}{!}{\begin{tikzpicture}
[scale=2,>=stealth,point/.style={draw,circle,fill = black,inner sep=0.5pt},]

\node (E) at (1.5, 1.5)[point,label=above :$E$] {};
\node (F) at (-1.5, 1.5)[point,label=above :$F$] {};
\node (A) at (0, 3)[point,label=above :$A$]{};
\node (B) at (-3, 0)[point,label=below left:$B$]{};
\node (C) at (3, 0)[point,label=below right:$C$]{};
\node (a) at (0,0)[point,label=below :$a$] {};
\node (O) at (0,1)[point,label=below :$O$] {};

\draw (a)--(B);
\draw (B)--(A);
\draw (A)--(C);
\draw (C)--(a);
\draw (B)--(E);
\draw (C)--(F);
\draw (E)--(F);
\node [above] at (-1.7,1.7) {$c$};
\node [above] at (1.7,1.7) {$b$};
\node [above] at (-2.5,.75) {$c/2$};
\node [above] at (-1,2.1) {$c/2$};
\node [above] at (2.5,.75) {$b/2$};
\node [above] at (1,2.1) {$b/2$};
\node [below] at (1,1.3) {$p$};
\node [below] at (-1,1.3) {$q$};
\node [above] at (-.7,.5) {$k_{p}$};
\node [above] at (.7,.5) {$k_{q}$};

\tkzMarkAngle[size=.3](F,O,B);
\tkzMarkAngle[size=.3](C,O,E);
\draw (-0.5,1) node{$\theta$};
\draw (0.5,1) node{$\theta$};

\end{tikzpicture}}
%	\end{center}
%	\caption{$O$ divides medians in the ratio $2:1$}
%	\label{ch2_median_ratio_val}	
%\end{figure}
%

%\proof Using the sine formula and \eqref{ch2_small_triang}, areas of some triangles in Fig. \ref{ch2_median_ratio_val} are listed in the following table
%\begin{center}
%\begin{tabular}{|c|c|}
%	\hline  Triangle& Area  \\ 
%	\hline  OFE & $\frac{1}{2}pq \sin \theta $ \\ 
%	\hline  BOF &  $\frac{k}{2}pq \sin \theta$\\ 
%	\hline  COE & $\frac{k}{2}pq \sin \theta$ \\ 
%	\hline  BOC&  $\frac{k^2}{2}pq \sin \theta$\\
%	\hline  BOC&  $\frac{1}{2} ar \brak{\Delta ABC}$\\  
%	\hline 
%\end{tabular} 
%\end{center}
%
%where we have used the fact that 
%
%\begin{equation}
%\sin \angle BOC = \sin \brak{180^{\degree}-\theta} = \sin \theta
%\end{equation}
%%
%Since $BE$ is the median
%\begin{align}
%\begin{split}
%ar\brak{\Delta BEC} &= \frac{1}{2} ar\brak{\Delta ABC}  
%\\
%& = ar\brak{\Delta BOC} + ar\brak{\Delta COE}\\
%ar\brak{\Delta BEA} &= \frac{1}{2} ar\brak{\Delta ABC}
%\\
%& = ar\brak{\Delta AFE} + ar\brak{\Delta BOF} +  ar\brak{\Delta FOE}  
%\end{split}
%\end{align}
%
%Substituting the respective values from the above table,
%\begin{align}
%\begin{split}
%\frac{1}{2} ar\brak{\Delta ABC} &= \frac{k^2}{2}pq \sin \theta + \frac{k}{2}pq \sin \theta
%\\
%\frac{1}{2} ar\brak{\Delta ABC} &= \frac{k}{2}pq \sin \theta + \frac{1}{2}pq \sin \theta + \frac{1}{4} ar\brak{\Delta ABC}
%\end{split}
%\end{align}
%%
%Simplifying the above,
%%
%\begin{align}
%k(k+1) = 2(k+1)
%\end{align}
%%
%Since $k \neq -1$, $k = 2$ and the proof is complete.
%
%
%\begin{problem}
%	In Fig. \ref{ch2_median_3}, the line $AO$ cuts $EF$ at $G$ and is extended to meet the side $BC$ at $D$.  Show that 
%	%
%	\begin{equation}
%	\frac{OA}{OD} = 2.
%	\end{equation}
%	%  
%\end{problem}
%\begin{figure}[!ht]
%	\begin{center}
%		
%		%\includegraphics[width=\columnwidth]{./figs/ch2_median_3}
%		%\vspace*{-10cm}
%		\resizebox{\columnwidth}{!}{\begin{tikzpicture}
[scale=2,>=stealth,point/.style={draw,circle,fill = black,inner sep=0.5pt},]

\node (E) at (1.5, 1.5)[point,label=above :$E$] {};
\node (F) at (-1.5, 1.5)[point,label=above :$F$] {};
\node (A) at (0, 3)[point,label=above :$A$]{};
\node (B) at (-3, 0)[point,label=below left:$B$]{};
\node (C) at (3, 0)[point,label=below right:$C$]{};
\node (O) at (0,1)[point,label=below right :$O$] {};
\node (D) at (0,0)[point,label=below :$D$] {};

\draw (B)--(A);
\draw (A)--(C);
\draw (C)--(B);
\draw (B)--(E);
\draw (C)--(F);
\draw (E)--(F);
\draw (A)--(D);

\node [above] at (-1.7,1.7) {$c$};
\node [above] at (1.7,1.7) {$b$};
\node [above] at (-2.5,.75) {$c/2$};
\node [above] at (-1,2.1) {$c/2$};
\node [above] at (2.5,.75) {$b/2$};
\node [above] at (1,2.1) {$b/2$};
\node [below] at (1,1.3) {$p$};
\node [below] at (-1,1.3) {$q$};
\node [above] at (-.7,.5) {$2p$};
\node [above] at (.7,.5) {$2q$};
\node [above] at (0,-.5) {$a$};
\node [above] at (0,1.9) [point,label=above right :$r$] {};
\node (G) at (0,1.5)[point,label=above right :$G$] {};

\tkzMarkAngle[size=.3](F,O,B);
\tkzMarkAngle[size=.3](C,O,E);
\draw (-0.5,1) node{$\theta$};
\draw (0.5,1) node{$\theta$};

\end{tikzpicture}}
%	\end{center}
%	\caption{Similar Triangles and Median}
%	\label{ch2_median_3}	
%\end{figure}
%
%\proof Since $EF || BC$,
%%
%\begin{equation}
%\Delta AGE \sim \Delta ADC \Rightarrow AG = GD = \frac{AD}{2}
%\end{equation}
%%
%Also,
%\begin{align}
%\Delta OGE \sim \Delta ODB \Rightarrow OD = 2r
%\end{align}
%Thus, we have the following relations
%\begin{align}
%GD &= OG + OD  = r + 2r = 3r\\
%AG &= GD = 3r\\
%OA &= OG + AG = r + 3r = 4r
%\end{align}
%Hence
%%
%\begin{equation}
%\frac{OA}{OD} = \frac{4r}{2r} = 2
%\end{equation}
%%
%
\begin{problem}
	In Fig. \ref{ch2_median_final},$BE$ is a median of $\Delta ABC$ and $\frac{OB}{OE} = 2$.  Show that AD is also a median. 
\end{problem}
%
\begin{figure}[!ht]
	\begin{center}
		
		%\includegraphics[width=\columnwidth]{./figs/ch2_median_final}
		%\vspace*{-10cm}
		\resizebox{\columnwidth}{!}{\begin{tikzpicture}
[scale=2,>=stealth,point/.style={draw,circle,fill = black,inner sep=0.5pt},]

\node (D) at (1, 0)[point,label=below :$D$] {};
\node (A) at (0, 3)[point,label=above :$A$]{};
\node (B) at (-3, 0)[point,label=below left:$B$]{};
\node (C) at (3, 0)[point,label=below right:$C$]{};
\node (E) at (1.5, 1.5)[point,label=above right:$E$]{};
\node (O) at (.6, 1.2)[point,label=above right:$O$]{};

\draw (D)--(B);
\draw (B)--(A);
\draw (A)--(C);
\draw (C)--(D);
\draw (D)--(A);
\draw (D)--(E);
\draw (E)--(B);
%\tkzMarkRightAngle[size=.2](A,D,C)

\node [below] at (0,-0.1) {$a$};
\node [above] at (-1.5,1.5){$c$};
\node [above] at (1.7,1.7){$b$};
\node [above] at (.7,.5){$x$};
\node [above] at (.19,2){$2x$};
\node [above] at (-.7,.5){$2p$};
\node [above] at (1.2,1.2){$p$};
\node [above] at (.9,2.2){$b/2$};
\node [above] at (2.3,0.8){$b/2$};

\tkzMarkAngle[size=.2](A,O,B);
\draw (0.3,1.3) node{$\theta$};
\tkzMarkAngle[size=.2](D,O,E);
\draw (0.9,1.1) node{$\theta$};

\end{tikzpicture}}
	\end{center}
	\caption{Medians meet at a point}
	\label{ch2_median_final}	
\end{figure}

\proof Since BE is a median, the areas of  triangles  $\triangle BEC$ and $\triangle BEA$ are equal. Hence,
\begin{align}
\frac{1}{2}BE.BC\sin \theta_1 &= \frac{1}{2}BE.AB\sin \theta_2
\\
\implies a\sin \theta_1 &= c\sin \theta_2
\label{eq:med_conv_area}
\end{align}
%
Using the sine formula in $\triangle OBD$ and $\triangle BOA$,
\begin{align}
\label{eq:med_conv_sin1}
\frac{q}{\sin \theta_1} &= \frac{x}{\sin \alpha} 
\\
\frac{c}{\sin \alpha} &= \frac{kq}{\sin \theta_2}
\label{eq:med_conv_sin2}
\end{align}
Mutiplying \eqref{eq:med_conv_area},\eqref{eq:med_conv_sin1}
and \eqref{eq:med_conv_sin2},
%
we obtain
\begin{align}
x = \frac{a}{k}
\label{eq:med_conv_ak}
\end{align}
after simplification.
Let the area of $\triangle ABC = \Delta$. Then, in terms of area,
\begin{align}
\triangle ABD &= \triangle OBD + \triangle OBA
\\
\implies \frac{1}{2}(c)(x)\sin B &=\frac{1}{2}(2p)(q)\sin \alpha 
\\
&\quad + \frac{1}{2}(2p)(kq)\sin \alpha
\\
\implies k\brak{k+1}pq \sin \alpha  &= \frac{1}{2}ca \sin B = \Delta
\label{eq:med_conv_del1}
\end{align}
%
Similarly,
\begin{align}
\triangle ABE &= \triangle AOB + \triangle AOE
\\
\implies \frac{\Delta}{2} &=\frac{1}{2}(2p)(kq)\sin \alpha 
\\
&\quad + \frac{1}{2}(p)(kq)\sin \alpha
\\
\implies 3kpq \sin \alpha  &=  \Delta
\label{eq:med_conv_del2}
\end{align}
%
after simplification. From \eqref{eq:med_conv_del1}
and \eqref{eq:med_conv_del2},
%
\begin{align}
k\brak{k+1}pq \sin \alpha &=  3kpq \sin \alpha  
\\
\implies k = 2.
\end{align}
Thus, 
\begin{align}
BD &= DC
\\
\frac{OA}{OD} &= 2.
\end{align}
Thus, $AD$ is a median.
\\
{\em Conculusion:} The medians of a triangle meet at a point.
